% state/part_intro.tex

% This file is included immediately after:
%   \part{Railway State Model}
%
% It is intentionally unnumbered. The numbered chapters of this part
% begin with pose2.tex.

\chapter*{Railway State Model}
\addcontentsline{toc}{chapter}{Railway State Model}

\begin{remark}
This part introduces the railway state model underlying the AIM
framework.
Instead of treating the track merely as a geometric curve, AIM describes
railway alignment through states generated from intrinsic geometry,
profile, cant, frames, and metric interpretation.
\end{remark}

\begin{remark}
The first state layer is \(\mathrm{pose2}\), the intrinsic planar
reconstruction state generated from curvature.
It carries the initial-value structure required to reconstruct position
and direction along the alignment.
\end{remark}

\begin{remark}
The second state layer is \(\mathrm{pose3}\), the spatial railway state
obtained by extending pose2 through profile, cant, and frame
interpretation.
This state describes the railway, not the vehicle.
\end{remark}

\begin{remark}
The purpose of this part is to define the variables, frames, and
transformations required to connect intrinsic alignment modelling with
runtime railway interpretation.
\end{remark}

\begin{principle}[Railway-state principle]
\label{princ:railway-state}
Within AIM, railway alignment is interpreted through state evolution
rather than through coordinate geometry alone.
\end{principle}

\begin{remark}
This part prepares the transition from geometric reconstruction toward
dynamic railway state interpretation, vehicle-relative states, and
runtime evaluation.
\end{remark}
